\section{万有引力定律}
\subsection{万有引力定律的数学形式}
在高中我们就学过万有引力定律了，如果有两个质点，它们的质量分别是$m_1$和$m_2$，它们之间的距离为$r$，那么它们之间的引力大小为
\[
F = G \frac{m_1 m_2}{r^2}
\]
这样的式子是一个标量式，如果想研究复杂的问题，矢量的形式可以免去几何问题，更加容易。我们用$\mathbf{F}_{12}$来表示$m_1$对$m_2$的引力，用$\hat{\mathbf{r}}_{12}$来表示由$m_1$指向$m_2$的单位矢量，由于引力是吸引力，所以方向与$\hat{\mathbf{r}}_{12}$相反。所以万有引力定律为
\[
\mathbf{F}_{12} = -G \frac{m_1 m_2}{r^2} \hat{\mathbf{r}}_{12}
\]
或利用$\hat{\mathbf{r}}_{12}= \mathbf{r}_{12}/r$，上面的式子可以写成
\[
\mathbf{F}_{12} = -G \frac{m_1 m_2}{r^3} \mathbf{r}_{12}
\]
这个式子是一种更常见的形式。

我们接下来将展示矢量形式的好处。
\subsection{质点组对质点的引力}
设空间中有 $N$ 个质点，第 $i$ 个质点的质量为 $m_i$，位置矢量为 $\mathbf{r}_i$。另有一质量为 $M$ 的质点位于 $\mathbf{r}$ 处。接下来我们尝试利用前面的定律求这 $N$ 个质点对 $M$ 的万有引力的合力。

先计算单个质点 $m_i$ 对 $M$ 的引力。令 $\mathbf{R}_{i} = \mathbf{r} - \mathbf{r}_i$ 为从 $m_i$ 指向 $M$ 的位矢，则引力大小为 $G \frac{m_i M}{|\mathbf{R}_{i}|^2}$，方向沿 $\hat{\mathbf{R}}_i = \mathbf{R}_i / |\mathbf{R}_i|$。故
\[
\mathbf{F}_i = -G \frac{m_i M}{|\mathbf{R}_i|^3} \mathbf{R}_i.
\]
所有质点对 $M$ 的总引力为各 $\mathbf{F}_i$ 的矢量和：
\[
\mathbf{F} = \sum_{i=1}^{N} \mathbf{F}_i = G M \sum_{i=1}^{N} m_i \frac{\mathbf{r}_i - \mathbf{r}}{|\mathbf{r} - \mathbf{r}_i|^3}.
\]
假如我们坚持用标量式，恐怕难以应对如此多的质点。

如果我们想研究一个在空间中质量连续分布的物体，那么就无法使用上面的结果了。但是在微积分里，我们通常可以把一个连续的东西，如计算函数曲线与坐标轴围成的面积，拆成若干可以计算的、离散的微元求和。

我们可以把一个连续的物体分割成无数个极小的质量微元 $\mathrm{d}m$，每个微元的尺度趋于零，从而可以视为一个质点，位于 $\mathbf{r}'$ 处。该微元对 $M$ 的引力为
\[
\mathrm{d}\mathbf{F} = -G M \, \mathrm{d}m \, \frac{\mathbf{r} - \mathbf{r}'}{|\mathbf{r} - \mathbf{r}'|^3}.
\]
引入质量密度 $\rho(\mathbf{r}')$，则 $\mathrm{d}m = \rho(\mathbf{r}') \, \mathrm{d}V'$，将所有微元的贡献叠加（即积分），便得到整个连续体对 $M$ 的总引力
\[
\mathbf{F} = -G M \iiint_V \rho(\mathbf{r}') \frac{\mathbf{r} - \mathbf{r}'}{|\mathbf{r} - \mathbf{r}'|^3} \, \mathrm{d}V'.
\]
\subsection{万有引力的性质}
万有引力的方向和大小赋予了其极其独特的性质。
\subsubsection{有心力}
若力 $\mathbf{F}$ 的方向始终通过一个固定点（称为力心），且大小仅依赖于质点至力心的距离 $r$，即 $\mathbf{F} = F(r) \hat{\mathbf{r}}$，则称 $\mathbf{F}$ 为有心力。万有引力就是一种有心力，$m_1$和$m_2$受到的引力始终指向两者的质心，也仅和两者之间的距离有关。

有心力具有许多共有的性质，包括运动将在一个固定的平面内（开普勒第一定律）、角动量守恒（开普勒第二定律）、具有对应的势能（引力势能）等。我们这里不做过多展开，如果对能量比较了解，可以试着求一下引力势能。
\subsubsection{对称球体的引力}
引力的平方反比的特性也赋予了它很独特的魅力。正如我们高中学到的库仑力一样，引力和库仑力是平方反比力。平方反比律导致球对称的球壳内部任意一点的物体将不受这个球壳的引力，对称球或球壳外部的质点受到的引力\textbf{等效于集中质量于球心处的质点对于它的引力}。在静电场中也有类似性质。这一性质的证明比较复杂，可以参考电磁学的教材来类比。

例如我们将地球视为均匀球体，质量为$M$，距离地球$r$处质量为m的卫星受到地球的引力大小为
\[
F = G \frac{Mm}{r^2}
\]
这个结论可以通过三重积分得到，也可以利用上面的性质得到。
\subsection{引力质量与惯性质量}
一个物体可能具有多种与相互作用有关的性质，比如它的电荷量决定了它受到的库仑相互作用的大小、\textbf{引力质量}决定了它受到引力的大小、\textbf{惯性质量}决定了它的惯性的大小。有趣的是，人们利用实验发现，引力质量等于惯性质量，这说明与库仑力和电荷量不同，引力和引力质量似乎更特殊一些。

如果引力质量严格等于惯性质量，那么一个物体在引力场中的牛顿第二定律就可以写为
\[
m \mathbf{a} = -G \frac{M m}{r^2} \hat{\mathbf{r}}
\]
左侧和右侧的$m$分别是惯性质量和引力质量。两边消去$m$，则方程可以简化为
\[
\mathbf{a} = -G \frac{M}{r^2} \hat{\mathbf{r}}
\]
这个微分方程的解将会和质量不再相关，物体在引力场中的运动与其质量无关。
\subsection{万有引力定律的应用}
\subsubsection{重力加速度}
在一个质量为$M$的质点周围，与新放入的研究对象的质量$m$无关的加速度为
\[
\mathbf{a} = -G \frac{M}{r^2} \hat{\mathbf{r}}
\]
只需要用$m$乘上这个加速度，我们就能很容易地获得该物体所受的力。这个加速度是一个关于空间坐标的矢量函数（即矢量场）。

如果忽略地球自转，并认为地球是一个半径为$R$的球体，我们来计算地球表面附近这个加速度的大小。

首先我们先假设有一个质量为$m$的物体作为研究对象，它处于地球表面附近，那么根据前面我们提到的性质，我们可以把地球视为一个质量集中于地心的质点。这样我们可以很容易地求出我们的研究对象所受的力的大小
\[
F = G \frac{Mm}{r^2}
\]
把这个力的大小除以我们研究对象的质量，就能得到物体在地球表面附近的加速度，即重力加速度，大小为
\[
g =\frac{F}{m}= G \frac{M}{R^2}
\]

值得注意的是，上面的讨论仅限于忽略地球自转并把地球当作球体的情况，而实际情况要复杂得多。在实际情况下，需要考虑惯性离心力和科里奥利力等，重力和引力之间存在微小的差异。
\subsubsection{第一宇宙速度}
已知地球质量为M，半径为R，我们来求解在地球表面附近绕地球做圆周运动的卫星的速度$v$。

此问题我们在高中就已经学过了，但我们可以重新审视这个问题。设卫星质量为$m$，在自然坐标系中$\mathbf e_n$指向地心，牛顿第二定律为
\[
m \mathbf a = G \frac{Mm}{R^2} \mathbf e_n
\]
写成分量的形式
\[
a_\tau = 0
\]
\[
a_n =G \frac{M}{R^2}
\]
再利用自然坐标系的结论
\[
a_n = \frac{v^2}{R}
\]
联立即可得到和高中阶段一致的结论
\[
v = \sqrt{\frac{GM}{R}}
\]
\subsubsection{一道经典的问题（选读）}
\noindent \textbf{例：}设地球是一个密度均匀的球体，其密度为$\rho$，重力加速度为$g$。A地正下方有一个半径为$a$的球形空洞，球心O距离地面的距离为$h$（$a<h$），求该地重力加速度$g'$。

\noindent \textbf{分析：}这个问题要求解的是一个完整球体挖去一个球体后剩余部分对于A点处质量为$m$的物体的引力$F'$，随后求出$g' = F'/m$，实际上可以存在两种不同的理解方式。

一种思路是，认为完整地球对于A的引力$mg$来源于两部分，一部分是挖去一部分的球体对A的引力$mg'$，另一部分是被挖去的球体对A的引力$m\Delta g$，而$mg$和$m\Delta g$因为是球体的引力所以特别好求解。做差即可求出$g'$。

另一种思路比较有趣，我们认为实际上被挖去一部分的球体是密度为 $\rho$ 的完整地球和密度为$-\rho$的小球体组合而成的。求出两者的引力求和即可。

\noindent \textbf{解：}以竖直向上为正方向。对于质量为 $m$ 的质点位于地面 A 点，完整地球对它的引力为  
\[
F_0 = -mg
\]  
空洞体积为 $\frac{4}{3}\pi a^3$，缺失的质量为  
\[
m_{\text{hole}} = \rho \cdot \frac{4}{3}\pi a^3.
\]  
地面点到空洞球心的距离为 $h$，故空洞区域原本的物质对质点的引力为  
\[
F_{\text{hole}} = -G \frac{m_{\text{hole}} \, m}{h^2} =- G \frac{\rho \cdot \frac{4}{3}\pi a^3 \, m}{h^2}.
\]  
实际引力为完整地球引力减去空洞部分引力 
\[
F' = F_0 - F_{\text{hole}} = - mg + G \frac{4\pi \rho a^3 m}{3h^2}.
\]  
同时
\[
F' = -mg'
\]
因此该点的重力加速度为  
\[
g' = - \frac{F'}{m} = g - G \frac{4\pi \rho a^3}{3h^2}.
\]  
此即所求重力加速度的表达式。